\documentclass[11pt]{article}
\usepackage{amsmath}
\usepackage{amssymb}
\usepackage{geometry}
\geometry{a4paper, margin=1in}

\title{Mathematical Formulation of the HeROsim Co-Simulation Environment}
\author{}
\date{}

\begin{document}

\maketitle

\section{The System State Formulation}
At any given discrete scheduling interval $t$, the state of the edge-cloud continuum is captured by HeROsim and mathematically projected as an attributed bipartite graph $\mathcal{G}_t = (\mathcal{U}_t, \mathcal{V}_t, \mathcal{E}_t)$. 
\begin{itemize}
    \item \textbf{The Demand Set}: $\mathcal{U}_t$ represents the incoming batch of $n$ serverless tasks.
    \item \textbf{The Supply Set}: $\mathcal{V}_t$ represents the $m$ available heterogeneous compute replicas.
    \item \textbf{The Adjacency}: $\mathcal{E}_t$ represents the valid, reachable topological pathways.
\end{itemize}

\section{The Action Space (Placement Matrix)}
The goal of the orchestrator is to formulate a routing policy. Mathematically, a complete placement decision (an action) evaluated by HeROsim is defined as a binary assignment matrix $Y \in \{0, 1\}^{n \times m}$.

For any specific task $u \in \mathcal{U}_t$ and replica $v \in \mathcal{V}_t$:
\begin{equation}
y_{uv} = 
\begin{cases} 
1, & \text{if task } u \text{ is routed to replica } v \\ 
0, & \text{otherwise} 
\end{cases}
\end{equation}

Because the system explicitly allows for non-unique placements to exploit warm-container queues, the strictly enforced constraint on this matrix is simply that every task must be assigned to exactly one replica:
\begin{equation}
\sum_{v \in \mathcal{V}_t} y_{uv} = 1 \quad \forall u \in \mathcal{U}_t
\end{equation}

\section{The HeROsim Cost Function (Deconstructing RTT)}
The core scientific value of HeROsim is its ability to accurately calculate the physical cost of any assignment matrix $Y$. The optimal placement minimizes the total Round-Trip Time (RTT), which is a composite metric. 

We define the discrete-event simulation as a non-linear cost function $\mathcal{F}_{sim}$, such that the empirical RTT of a specific task $u$ executing on replica $v$ is modeled as:
\begin{equation}
RTT_{uv} = \tau_{net}(u,v) + \tau_{queue}(v, Y) + \big[\tau_{init}(v) \cdot \delta_{cold}(v)\big] + \tau_{exec}(u,v)
\end{equation}

Where:
\begin{itemize}
    \item $\tau_{net}(u,v)$: The topological network transmission and routing latency.
    \item $\tau_{queue}(v, Y)$: The dynamic queuing delay at the replica. Crucially, this depends on the \textit{entire} assignment matrix $Y$ (multi-tenant contention).
    \item $\tau_{init}(v)$: The hardware-specific container initialization time (the cold start penalty).
    \item $\delta_{cold}(v) \in \{0,1\}$: A binary state indicator from the replica's feature vector $x_v$, where $1$ denotes a cold state requiring initialization.
    \item $\tau_{exec}(u,v)$: The hardware-specific execution variance (e.g., CPU vs. DLA execution times).
\end{itemize}

The total objective cost for a given batch simulated by HeROsim is therefore:
\begin{equation}
RTT_{total}(Y) = \sum_{u \in \mathcal{U}_t} \sum_{v \in \mathcal{V}_t} y_{uv} \cdot RTT_{uv}
\end{equation}

\section{The Omniscient Oracle and Regret Generation}
To train the GNN, HeROsim acts as an omniscient oracle by performing an exhaustive enumeration of the combinatorial search space. Let $\Omega$ be the set of all valid placement matrices. 

During the offline generation phase, HeROsim calculates the absolute optimal ground-truth placement $Y_{best}$ and its resulting minimum latency $RTT_{best}$:
\begin{equation}
Y_{best} = \underset{Y \in \Omega}{\mathrm{argmin}} \ RTT_{total}(Y)
\end{equation}
\begin{equation}
RTT_{best} = RTT_{total}(Y_{best})
\end{equation}

Simultaneously, the simulator generates the comprehensive lookup table $\mathcal{L}$ for the batch, mapping every suboptimal alternative assignment $Y_{alt}$ to its corresponding penalty:
\begin{equation}
\Delta RTT(Y_{alt}) = RTT_{total}(Y_{alt}) - RTT_{best}
\end{equation}

This explicitly defined $\Delta RTT$ is the exact empirical regret that dynamically scales the proposed pairwise margin loss $\mathcal{L}_{regret}$ during the GNN training phase.

\end{document}